\section{Life Cycle of a Login, a Message and a File}\label{sec:lifecycles}

\subsection{Joining the network}
\Cref{fig:seq-login} follows the initialisation phase. When a peer daemon starts it binds its three
sockets and registers with \code{/submit-info}; the tracker answers with a peer token and the current
list, and the daemon repeats the call every \SI{10}{s} as a heartbeat. When the user signs in, the
daemon forwards the credentials to the tracker, which checks the password and records the user on
that peer; only then does the daemon start a cookie session for the browser. Discovery is
\code{/get-list}. Connection set-up is a \code{CONNECT}/\code{CONNECT_OK} exchange; both sides mark
each other as connected, so one click connects a pair.

\fig[fig:seq-login]{seq_login.pdf}{Registration, login, discovery and connection set-up}{0.92\textwidth}

\subsection{A chat message}
\Cref{fig:seq-message} shows a broadcast in \code{\#general} with one unreachable peer. The daemon
builds the frame, encrypts the text, stores its own copy, and sends the frame to all connected peers
in parallel (a pool of 16 threads, \SI{2}{s} connect timeout). Each receiver reassembles the frame,
decrypts it, drops it if the \code{message_id} is already known, and appends it to its history with a
\emph{local} sequence number. Peers that could not be reached are passed to the tracker's
\code{/relay}, which forwards the frame unchanged --- it has no chat key and cannot read it. The
answer to the browser reports which peers got the message directly, which via the relay and which not
at all.

\fig[fig:seq-message]{seq_message.pdf}{Life cycle of a broadcast message, with the relay fallback}{\textwidth}

A direct message follows the same path with \code{to} set to the recipient; the recipient files it
under the conversation \code{@<sender peer id>} and every other peer would discard it. The browser
polls \code{/api/messages?after=<seq>} every \SI{1.5}{s}. Using the receiver's own sequence number as
the cursor matters: an earlier version used the sender's timestamp, so a message from a peer whose
clock was behind was never shown (\cref{sec:challenges}).

\pysnippet[lst:fanout]{fan_out.py}{Parallel delivery with the relay fallback}

\subsection{A file}
\Cref{fig:seq-file} follows a file. The browser uploads it to its own daemon, which answers
\code{202 Accepted} with a transfer id and sends the file in a worker thread; the browser polls
\code{/api/transfers} to draw the progress bar. On the wire, \code{START} is acknowledged before any
data is sent. The sender then keeps up to \code{cwnd} chunks in flight. In the figure chunk~1 is lost;
the three acknowledgements of chunks 2--4 all carry \code{cum = 1}, which tells the sender that chunk
1 is missing, and it is resent at once instead of after a timeout. The last acknowledgement carries
\code{complete = true} only after the receiver has checked the SHA-256 of the whole file and renamed
the temporary \code{.part} file.

\fig[fig:seq-file]{seq_file.pdf}{Life cycle of a file transfer with one lost chunk}{0.78\textwidth}
